% !TEX root = 00_instructivo.tex
\chapter{Transductores acústicos: hidrófono con piezoeléctrico y amplificador MAX4466}

\section*{Objetivos}
\begin{itemize}
    \item Construir un hidrófono usando un sensor piezoeléctrico impermeabilizado.
    \item Acondicionar la señal con el amplificador MAX4466.
    \item Clasificar sonidos submarinos (golpes, motores, burbujas) mediante umbrales de amplitud.
\end{itemize}

\section*{Materiales y equipo}
\begin{itemize}
    \item Sensor piezoeléctrico analógico Keyestudio (o disco piezo con cables).
    \item Módulo amplificador MAX4466 (ganancia ajustable).
    \item Arduino Uno.
    \item Mini protoboard de 9 canales.
    \item Recipiente con agua (pecera).
    \item Fuentes de sonido: motor pequeño, objeto metálico, manguera para burbujas.
    \item Silicona o resina epóxica para impermeabilizar.
\end{itemize}

\section*{Instrucciones}
\begin{enumerate}
    \item Cubra el piezoeléctrico con una capa delgada de silicona, dejando el conector expuesto. Deje secar.
    \item Conecte el piezo a la entrada del MAX4466 (IN y GND). La salida del MAX4466 va a un pin analógico del Arduino (A0).
    \item Ajuste la ganancia del MAX4466 para que el ruido de fondo en agua tranquila sea ~1.5V DC (offset).
    \item Genere sonidos: golpee la pared del recipiente, active el motor sumergido, produzca burbujas.
    \item Programe umbrales: si la amplitud pico supera un valor V1 encienda LED rojo, si está entre V0 y V1 LED azul.
    \item Registre 10 segundos de señal para cada fuente y calcule el nivel RMS relativo.
\end{enumerate}

\section*{Esquemático}
\begin{figure}[htbp]
\centering
\begin{circuitikz}[scale=0.7]
\draw (0,0) node[above]{Arduino};
% MAX4466
\draw (2,0) node[op amp, scale=0.8] (max) {};
\draw (max.out) -- (3,0) node[anchor=west]{A0};
\draw (1,-0.5) node[anchor=east]{5V} -- (max.V+);
\draw (1,-1) node[anchor=east]{GND} -- (max.V-);
% Piezo conectado a entrada del MAX4466
\draw (0,0.5) node[pz] (piezo) {};
\draw (piezo.+) -- (max.-);
\draw (piezo.-) -- (max.+);
\end{circuitikz}
\caption{El piezo se conecta a la entrada diferencial del MAX4466. La salida va a A0. Alimentación 5V y GND.}
\end{figure}

\section*{Preguntas de análisis}
\begin{enumerate}
    \item ¿Por qué se necesita un amplificador de ganancia ajustable? ¿Qué pasa si la ganancia es muy alta?
    \item ¿Cómo afecta la impermeabilización a la sensibilidad del piezoeléctrico?
    \item Compare la respuesta en frecuencia de un hidrófono piezoeléctrico vs un micrófono convencional.
    \item Proponga un método para clasificar sonidos no solo por amplitud sino por contenido espectral (FFT).
\end{enumerate}

\section*{Bibliografía}
\begin{itemize}
    \item MAX4466 Datasheet. Maxim Integrated.
    \item “Piezoelectric hydrophone” – Application note, Measurement Specialties.
    \item Keyestudio Analog Piezoelectric Ceramic Vibration Sensor – User manual.
\end{itemize}
